\cleardoublepage
\chapter{KESIMPULAN DAN SARAN}
\label{chap:penutup}

Bab ini merangkum temuan-temuan utama dari hasil perancangan, implementasi, dan pengujian sistem \textit{Smart Waste Bin}, serta menguraikan rekomendasi praktis dan akademis untuk pengembangan penelitian maupun penyempurnaan perangkat di masa mendatang.

% ============================================================
\section{Kesimpulan}
\label{sec:kesimpulan}
% ============================================================

Penelitian ini telah berhasil mengembangkan dan mengimplementasikan sistem \textit{Smart Waste Bin} berbasis \textit{Internet of Things} (IoT) untuk keperluan otomatisasi pengumpulan data dan analisis pola pembuangan sampah di Gedung Labtek VIII Institut Teknologi Bandung, khususnya pada Laboratorium \textit{Internet of Things} Lantai 4. Selama periode observasi 37 hari, sistem berhasil mengumpulkan 816 \textit{records} data dari tiga unit perangkat sensor dengan tingkat keberhasilan pengiriman mencapai 91,9\%, dengan 805 \textit{records} valid setelah \textit{filtering} data kotor. Berdasarkan hasil evaluasi menyeluruh pada Bab VI, seluruh rumusan masalah dalam penelitian ini telah berhasil dijawab dengan rincian sebagai berikut.

\begin{enumerate}
  \item \textbf{Rancangan dan Implementasi Sistem (Menjawab RM-1)}

  Penelitian ini berhasil merancang sistem \textit{Smart Waste Bin} terintegrasi dengan arsitektur IoT 4-\textit{layer}. Arsitektur ini mencakup \textit{Perception Layer} (ESP32, VL53L0X, INA219, SIM800L) untuk akuisisi data fisik, \textit{Network Layer} (MQTT) untuk transmisi nirkabel, \textit{Middleware Layer} (Node-RED) untuk pemrosesan aliran data, dan \textit{Application Layer} (Supabase, Next.js) untuk penyimpanan dan visualisasi antarmuka. Seluruh komponen terbukti berfungsi sesuai dengan desain, di mana 10 kebutuhan fungsional dan 7 dari 8 kebutuhan nonfungsional telah terpenuhi.

  \item \textbf{Identifikasi Pola Pembuangan Sampah (Menjawab RM-2)}

  Melalui analisis data historis dan analisis delta, sistem berhasil mengidentifikasi empat pola temporal pembuangan sampah di lingkungan kampus yang telah divalidasi melalui observasi lapangan, yaitu:

  \begin{enumerate}[a.]
    \item \textit{Peak hour} (jam sibuk) secara konsisten pada pukul 14:00 sampai 17:00, ditunjukkan oleh 8 \textit{event} penambahan (29,6\% dari total 27 \textit{event}). Pola ini konsisten dengan perilaku membuang sampah setelah aktivitas makan siang.

    \item Seluruh \textit{event} penambahan sampah terjadi pada jam aktif hari kerja (07:00 sampai 19:00), tanpa satu pun penambahan yang terdeteksi pada malam hari maupun akhir pekan.

    \item Aktivitas pembuangan sampah turun drastis pada hari libur nasional, di mana tidak terdeteksi adanya \textit{event} penambahan pada 25 sampai 26 Desember 2025 dan 1 Januari 2026. Nilai \textit{fillLevel} tinggi pada hari libur merupakan efek akumulasi kumulatif, bukan cerminan aktivitas pembuangan.

    \item Dominasi sampah anorganik yang menyumbang 13 dari 27 \textit{event} penambahan (48,1\%) dengan total delta positif sebesar 150\%, konsisten dengan karakteristik lokasi sebagai area laboratorium yang menghasilkan banyak sampah kemasan sekali pakai.
  \end{enumerate}

  \item \textbf{Peningkatan \textit{Firmware Smart Waste Bin} (Menjawab RM-3)}

  Perangkat bawaan dari pihak ketiga telah berhasil dioptimasi melalui lima peningkatan \textit{firmware} agar sesuai dengan kebutuhan lingkungan kampus, meliputi perubahan interval pengambilan data, implementasi \textit{median filter}, mekanisme \textit{auto-recovery}, \textit{warm-up readings}, dan GPIO \textit{isolation}. Empat dari lima peningkatan terbukti efektif dalam mengatasi masalah yang ditargetkan. GPIO \textit{isolation} tidak menunjukkan perubahan signifikan pada konsumsi daya karena sumber konsumsi utama berasal dari modul GSM yang tidak terpengaruh oleh kondisi GPIO. Peningkatan yang terbukti efektif dapat direkomendasikan untuk diadopsi pada produk \textit{Smart Waste Bin} standar vendor, khususnya untuk \textit{deployment} di lingkungan dengan karakteristik aktivitas yang terkonsentrasi pada jam-jam tertentu.
\end{enumerate}

% ============================================================
\section{Saran}
\label{sec:saran}
% ============================================================

Penelitian ini masih memiliki sejumlah keterbatasan yang membuka peluang untuk pengembangan lebih lanjut. Berdasarkan temuan dan kendala yang dihadapi selama pelaksanaan penelitian, saran dibagi menjadi dua kelompok, yaitu saran untuk penelitian selanjutnya dan saran untuk penyedia teknologi \textit{Smart Waste Bin}.

\subsection{Saran untuk Penelitian Selanjutnya}

Bagi peneliti yang ingin mengembangkan penelitian serupa, terdapat beberapa aspek yang dapat ditingkatkan sebagai berikut.

\begin{enumerate}
  \item Penggunaan tempat sampah berukuran lebih kecil disarankan agar resolusi pengukuran \textit{fillLevel} menjadi lebih tinggi. Pada penelitian ini, pengukuran terbatas pada kelipatan 10\% akibat penggunaan tipe data integer pada \textit{firmware}.

  \item Periode pengumpulan data perlu diperpanjang minimal satu semester untuk memperoleh pola yang lebih representatif. Periode yang lebih panjang juga memungkinkan validasi pola mingguan dengan tingkat keyakinan yang lebih tinggi.

  \item Penambahan sensor pada beberapa titik (\textit{multi-point}) dapat dipertimbangkan untuk mendeteksi sampah yang menumpuk di pinggir tempat sampah. Kondisi ini tidak terdeteksi pada penelitian ini karena menggunakan konfigurasi sensor \textit{single-point}.

  \item Program sosialisasi pemilahan sampah dapat dikombinasikan dengan penerapan sistem untuk menekan tingkat kesalahan klasifikasi sampah, sebagaimana ditemukan pada observasi lapangan dalam penelitian ini.
\end{enumerate}

\subsection{Saran untuk Penyedia Teknologi \textit{Smart Waste Bin}}

Bagi penyedia perangkat \textit{Smart Waste Bin}, beberapa rekomendasi berikut dapat dipertimbangkan untuk meningkatkan kualitas produk.

\begin{enumerate}
  \item Ditemukan variasi konsumsi arus hampir dua kali lipat antar unit dengan spesifikasi identik, yaitu 110 sampai 140 mA pada dua unit dan 210 sampai 235 mA pada satu unit. Investigasi terhadap kebocoran daya pada jalur komponen perlu dilakukan, termasuk penerapan \textit{quality control} dan \textit{burn-in testing} sebelum perangkat dipasarkan.

  \item Empat peningkatan \textit{firmware} yang terbukti efektif, yaitu perubahan interval, \textit{median filter}, \textit{auto-recovery}, dan \textit{warm-up readings}, dapat diadopsi ke produk standar untuk meningkatkan keandalan sistem, khususnya di lingkungan dengan karakteristik aktivitas yang dinamis.

  \item Dokumentasi teknis perlu dilengkapi dengan skema rangkaian dan \textit{source code firmware} agar mempermudah proses pemeliharaan dan \textit{troubleshooting} oleh pengguna.
\end{enumerate}